% ============================================================================
% Section 3: Stochastic Test Semantics
% ============================================================================

\section{Stochastic Test Semantics}
\label{sec:stochastic}

This section presents the formal foundation of \agentassay{}: a
stochastic test semantics that replaces binary pass/fail verdicts with
three-valued probabilistic outcomes backed by statistical guarantees.

\subsection{Preliminaries}
\label{sec:stochastic:prelim}

We begin by formalizing the objects under test.

\begin{definition}[Agent]
\label{def:agent}
An \emph{agent} is a tuple $A = (\pi, \mathcal{T}, \mu, \omega)$ where:
\begin{itemize}[leftmargin=*]
  \item $\pi$ is the agent's \emph{prompt} (system instructions, persona,
    goals),
  \item $\mathcal{T} = \{t_1, t_2, \ldots, t_k\}$ is the set of
    available \emph{tools},
  \item $\mu$ is the underlying \emph{language model} (including its
    version, temperature, and sampling parameters), and
  \item $\omega$ is the \emph{orchestration logic} (the control flow
    governing tool selection and multi-step execution).
\end{itemize}
\end{definition}

\begin{definition}[Agent Execution Trace]
\label{def:trace}
An \emph{execution trace} of agent $A$ on input $x$ is a sequence
$\tau = (s_1, s_2, \ldots, s_m)$ where each \emph{step} $s_i$ is a
tuple:
\[
  s_i = (a_i, t_i, o_i, c_i)
\]
with $a_i \in \mathcal{A}$ the \emph{action} (reason, call\_tool,
respond), $t_i \in \mathcal{T} \cup \{\bot\}$ the \emph{tool invoked}
(or $\bot$ if no tool call), $o_i$ the \emph{output} of the step, and
$c_i \in \mathbb{R}_{\geq 0}$ the \emph{cost} (API tokens, latency, monetary).
The \emph{final output} of a trace is $\mathrm{out}(\tau) = o_m$.
\end{definition}

Because agent $A$ is stochastic, repeated executions on the same input
$x$ produce a \emph{distribution} over traces. We write
$\tau \sim A(x)$ to denote a trace sampled from this distribution.

\begin{definition}[Evaluator]
\label{def:evaluator}
An \emph{evaluator} is a function $E: \mathcal{X} \times
\mathcal{O} \to \{0, 1\}$ that maps an input $x$ and an output $o$ to a
binary judgment: $E(x, o) = 1$ if the output is acceptable and
$E(x, o) = 0$ otherwise. An evaluator may be:
\begin{enumerate}[leftmargin=*]
  \item \emph{Deterministic}: a rule-based check (e.g., ``output
    contains keyword $k$''),
  \item \emph{Model-based}: an LLM judge with its own stochasticity, or
  \item \emph{Contract-based}: a behavioral contract from
    \agentassert{}~\citep{bhardwaj2026abc}.
\end{enumerate}
\end{definition}

\begin{remark}
When the evaluator is itself stochastic (model-based), the test
introduces a second source of randomness. Our framework handles this by
treating the composition of agent and evaluator as a single stochastic
process. We discuss this further in \cref{sec:discussion}.
\end{remark}

\subsection{Test Scenarios and the Test Triple}
\label{sec:stochastic:scenario}

\begin{definition}[Test Scenario]
\label{def:scenario}
A \emph{test scenario} is a tuple $S = (x, \mathcal{P}, E)$ where:
\begin{itemize}[leftmargin=*]
  \item $x \in \mathcal{X}$ is the \emph{input} to the agent,
  \item $\mathcal{P}$ is the set of \emph{expected properties} (natural
    language descriptions, formal assertions, or contract references),
  \item $E$ is the \emph{evaluator} (\cref{def:evaluator}).
\end{itemize}
\end{definition}

\begin{definition}[$(\alpha, \beta, n)$-Test Triple]
\label{def:test-triple}
A \emph{stochastic test} is parameterized by a triple
$T = (\alpha, \beta, n)$ where:
\begin{itemize}[leftmargin=*]
  \item $\alpha \in (0, 1)$ is the \emph{significance level} (upper
    bound on Type~I error probability---the probability of declaring
    \FAIL{} when the agent actually meets the threshold),
  \item $\beta \in (0, 1)$ is the \emph{Type~II error probability}
    (upper bound on the probability of declaring \PASS{} when the agent
    is actually below the threshold; $1 - \beta$ is the \emph{power}),
  \item $n \in \N_{\geq 1}$ is the \emph{number of trials} (independent
    executions of the agent on the test scenario).
\end{itemize}
\end{definition}

The three parameters are not independent. For a given effect size
$\delta$ (the minimum regression magnitude to detect), the required
sample size is:
\begin{equation}
\label{eq:sample-size}
  n^*(\alpha, \beta, \delta) =
    \left\lceil
      \frac{(z_{1-\alpha/2} + z_{1-\beta})^2 \cdot
            2\hat{p}(1-\hat{p})}{\delta^2}
    \right\rceil
\end{equation}
where $z_q$ denotes the $q$-quantile of the standard normal distribution
and $\hat{p}$ is the estimated pass rate under $H_0$.

\noindent\textit{Remark.} Equation~\eqref{eq:sample-size} uses the
two-sided quantile $z_{1-\alpha/2}$ for the single-version threshold
test. The regression test (\cref{eq:required-n}) uses the one-sided
quantile $z_{1-\alpha}$, matching the one-sided alternative
$H_1: p_c < p_b$.

\subsection{The Verdict Function}
\label{sec:stochastic:verdict}

\begin{definition}[Stochastic Verdict]
\label{def:verdict}
Given a test scenario $S$, a test triple $T = (\alpha, \beta, n)$, a
pass threshold $\theta \in (0, 1)$, and trial results
$\mathbf{r} = (r_1, r_2, \ldots, r_n)$ with $r_i = E(x, \mathrm{out}(\tau_i))$
for independently sampled traces $\tau_i \sim A(x)$, define the
\emph{observed pass rate}:
\begin{equation}
\label{eq:passrate}
  \hat{p} = \frac{1}{n} \sum_{i=1}^{n} r_i
\end{equation}
and the Wilson score confidence interval at level $1-\alpha$:
\begin{equation}
\label{eq:wilson}
  \CI_{1-\alpha}(\hat{p}, n) = \left[
    \frac{\hat{p} + \frac{z^2}{2n} - z\sqrt{\frac{\hat{p}(1-\hat{p})}{n}
    + \frac{z^2}{4n^2}}}{1 + \frac{z^2}{n}},\;
    \frac{\hat{p} + \frac{z^2}{2n} + z\sqrt{\frac{\hat{p}(1-\hat{p})}{n}
    + \frac{z^2}{4n^2}}}{1 + \frac{z^2}{n}}
  \right]
\end{equation}
where $z = z_{1-\alpha/2}$. Let
$\CI_{\text{lower}} = \CI_{1-\alpha}(\hat{p}, n)_{\text{lower}}$ and
$\CI_{\text{upper}} = \CI_{1-\alpha}(\hat{p}, n)_{\text{upper}}$. The
\emph{verdict function} is:
\begin{equation}
\label{eq:verdict}
  V(\mathbf{r}; \theta, \alpha) =
  \begin{cases}
    \PASS & \text{if } \CI_{\text{lower}} \geq \theta \\
    \FAIL & \text{if } \CI_{\text{upper}} < \theta \\
    \INCONC & \text{otherwise}
  \end{cases}
\end{equation}
\end{definition}

\begin{example}
\label{ex:verdict}
Consider an agent tested with $n = 50$ trials at threshold
$\theta = 0.85$ and significance $\alpha = 0.05$. If $\hat{p} = 0.90$
(45 of 50 pass), the Wilson 95\% CI is approximately $[0.789, 0.958]$.
Since $\CI_{\text{lower}} = 0.789 < 0.85 = \theta$, the verdict is
\INCONC{}---we cannot yet confirm the agent meets the threshold with 95\%
confidence. With $n = 100$ and $\hat{p} = 0.90$, the CI narrows to
approximately $[0.826, 0.946]$: still \INCONC{}. At $n = 200$ with
$\hat{p} = 0.90$, the CI becomes $[0.853, 0.935]$: now
$\CI_{\text{lower}} \geq \theta$, so the verdict is \PASS{}.
\end{example}

The three-valued verdict captures a fundamental truth: \emph{statistical
evidence can be insufficient}, and the framework should say so rather
than forcing a premature binary decision.

\begin{theorem}[Verdict Soundness]
\label{thm:soundness}
Let $p$ be the true (unknown) pass rate of agent $A$ on scenario $S$
with evaluator $E$. If $V(\mathbf{r}; \theta, \alpha) = \PASS$ under
test triple $(\alpha, \beta, n)$, then:
\begin{enumerate}[leftmargin=*]
  \item The probability that $V = \PASS$ when $p < \theta$ (false
    positive) satisfies $\Pr[V = \PASS \mid p < \theta] \leq \alpha$.
  \item The true pass rate satisfies $p \geq \theta - \varepsilon(n)$
    with probability at least $1 - \alpha$, where the tolerance
    $\varepsilon(n) = O(1/\sqrt{n})$ converges to zero.
\end{enumerate}
\end{theorem}

\begin{proof}[Proof sketch]
Part (1) follows from the coverage guarantee of the Wilson confidence
interval: $\Pr[p \in \CI_{1-\alpha}(\hat{p}, n)] \geq 1 - \alpha$. If
$p < \theta$, then for $V = \PASS$ to occur, we need
$\CI_{\text{lower}} \geq \theta > p$, which means $p$ falls below the
confidence interval. By the coverage guarantee, this happens with
probability at most $\alpha$.

Part (2) follows from inverting the confidence interval.
$V = \PASS$ implies $\CI_{\text{lower}} \geq \theta$, and since
$\CI_{\text{lower}} \leq p + O(1/\sqrt{n})$ with high probability, we
obtain $p \geq \theta - O(1/\sqrt{n})$.

The full proof, using the Clopper-Pearson exact interval for the formal
bound, is in \cref{app:proof:soundness}.
\end{proof}

\subsection{Regression Detection}
\label{sec:stochastic:regression}

\begin{definition}[Regression Verdict]
\label{def:regression}
Given baseline results $\mathbf{r}_b = (r_1^b, \ldots, r_{n_b}^b)$ with
pass rate $\hat{p}_b = \sum r_i^b / n_b$ and current results
$\mathbf{r}_c = (r_1^c, \ldots, r_{n_c}^c)$ with pass rate
$\hat{p}_c = \sum r_i^c / n_c$, define the \emph{regression verdict}:
\begin{equation}
\label{eq:regression-verdict}
  V_{\text{reg}}(\mathbf{r}_b, \mathbf{r}_c; \alpha, \beta, \delta) =
  \begin{cases}
    \FAIL & \text{if } p\text{-value}(H_0) < \alpha \text{ and }
      |\hat{p}_b - \hat{p}_c| \geq \delta \\
    \PASS & \text{if } p\text{-value}(H_0) \geq \alpha \text{ and }
      \text{power} \geq 1-\beta \\
    \INCONC & \text{otherwise}
  \end{cases}
\end{equation}
where $H_0: p_c \geq p_b$ (no regression), $H_1: p_c < p_b$
(regression occurred), the $p$-value is computed via Fisher's exact test
or the $Z$-test for proportions, and $\delta > 0$ is the minimum
\emph{clinically significant} effect size.
\end{definition}

The dual requirement of statistical significance ($p < \alpha$) \emph{and}
practical significance ($|\hat{p}_b - \hat{p}_c| \geq \delta$) follows
the recommendation of Arcuri and Briand~\citep{arcuri2011practical} and
prevents declaring regressions that are statistically detectable but
practically negligible.

\begin{theorem}[Regression Detection Power]
\label{thm:power}
Given a baseline pass rate $p_b$, a true regressed rate
$p_c = p_b - \delta$ for effect size $\delta > 0$, and test triple
$(\alpha, \beta, n)$: if the sample sizes satisfy
\begin{equation}
\label{eq:required-n}
  n_b, n_c \geq n^*(\alpha, \beta, \delta) =
  \left\lceil
    \frac{(z_{1-\alpha} + z_{1-\beta})^2 \cdot
          [p_b(1-p_b) + p_c(1-p_c)]}
         {\delta^2}
  \right\rceil
\end{equation}
where $p_b$ and $p_c$ are the baseline and current pass rates respectively, then the probability that the
regression verdict detects the regression satisfies:
\[
  \Pr[V_{\text{reg}} = \FAIL \mid p_c = p_b - \delta] \geq 1 - \beta.
\]
\end{theorem}

\begin{proof}[Proof sketch]
Under $H_1: p_c = p_b - \delta$, the test statistic
$Z = (\hat{p}_b - \hat{p}_c) / \text{SE}$ follows approximately
$\mathcal{N}(\delta / \text{SE}, 1)$. The critical value is
$z_{1-\alpha}$, and the probability of exceeding it under $H_1$ is
$\Pr[Z > z_{1-\alpha}] = \Phi(z_{1-\beta}) = 1-\beta$ when $n$ achieves
the required sample size. The full proof via the Neyman-Pearson lemma is
in \cref{app:proof:power}.
\end{proof}

\begin{definition}[Effect Size Measures]
\label{def:effect-size}
We quantify regression magnitude using three complementary measures:
\begin{enumerate}[leftmargin=*]
  \item \emph{Absolute difference}: $\Delta = \hat{p}_b - \hat{p}_c$
  \item \emph{Cohen's $h$}~\citep{cohen1988statistical}:
    $h = 2\arcsin(\sqrt{\hat{p}_b}) - 2\arcsin(\sqrt{\hat{p}_c})$
  \item \emph{Odds ratio}: $\text{OR} = \frac{\hat{p}_b / (1-\hat{p}_b)}
    {\hat{p}_c / (1-\hat{p}_c)}$
\end{enumerate}
Cohen's $h$ is the primary measure, with $|h| < 0.2$ indicating a small
effect, $0.2 \leq |h| < 0.5$ medium, and $|h| \geq 0.5$ large.
\end{definition}

\subsection{Sequential Probability Ratio Test}
\label{sec:stochastic:sprt}

Fixed-sample testing (\cref{sec:stochastic:verdict}) requires
pre-specifying $n$. For agents, where each trial incurs non-trivial cost,
an adaptive approach is preferable.

\begin{definition}[SPRT for Agent Testing]
\label{def:sprt}
Given hypotheses $H_0: p \geq \theta$ and $H_1: p \leq \theta - \delta$,
define the log-likelihood ratio after $k$ trials:
\begin{equation}
\label{eq:sprt-llr}
  \Lambda_k = \sum_{i=1}^{k} \log \frac{
    (\theta - \delta)^{r_i} (1 - \theta + \delta)^{1-r_i}
  }{
    \theta^{r_i} (1-\theta)^{1-r_i}
  }
\end{equation}
with boundaries $a = \log(\beta / (1-\alpha))$ and
$b = \log((1-\beta)/\alpha)$. The SPRT decision rule is:
\begin{equation}
\label{eq:sprt-rule}
  V_{\text{SPRT}} =
  \begin{cases}
    \PASS & \text{if } \Lambda_k \leq a \\
    \FAIL & \text{if } \Lambda_k \geq b \\
    \text{continue testing} & \text{if } a < \Lambda_k < b
  \end{cases}
\end{equation}
\end{definition}

\begin{proposition}[SPRT Efficiency]
\label{thm:sprt}
The SPRT (\cref{def:sprt}) satisfies the following classical properties
of Wald's sequential analysis~\citep{wald1947sequential}:
\begin{enumerate}[leftmargin=*]
  \item \textbf{Error control:} $\Pr[\text{accept } H_1 \mid H_0] \leq \alpha$
    and $\Pr[\text{accept } H_0 \mid H_1] \leq \beta$.
  \item \textbf{Sample efficiency:} The expected number of trials under
    $H_0$ and $H_1$ satisfies:
    \begin{align}
    \label{eq:sprt-h0}
      \E[N \mid H_0] &\approx
        \frac{(1-\alpha)\log\frac{\beta}{1-\alpha} +
              \alpha\log\frac{1-\beta}{\alpha}}
             {\theta\log\frac{\theta}{\theta-\delta} +
              (1-\theta)\log\frac{1-\theta}{1-\theta+\delta}} \\
    \label{eq:sprt-h1}
      \E[N \mid H_1] &\approx
        \frac{\beta\log\frac{\beta}{1-\alpha} +
              (1-\beta)\log\frac{1-\beta}{\alpha}}
             {(\theta-\delta)\log\frac{\theta-\delta}{\theta} +
              (1-\theta+\delta)\log\frac{1-\theta+\delta}{1-\theta}}
    \end{align}
  \item \textbf{Optimality:} Among all sequential tests with error
    probabilities at most $\alpha$ and $\beta$, the SPRT minimizes
    $\E[N \mid H_i]$ for $i \in \{0, 1\}$ (Wald-Wolfowitz theorem).
\end{enumerate}
\end{proposition}

\begin{proof}[Proof sketch]
Part (1) follows from Wald's fundamental identity for sequential tests.
Parts (2) and (3) follow from Wald's equation
$\E[\Lambda_N] = \E[N] \cdot \E[\Lambda_1]$ and the Wald-Wolfowitz
optimality theorem. See \cref{app:proof:sprt} for the full derivation.
\end{proof}

\begin{example}
\label{ex:sprt-savings}
Consider testing whether an agent maintains $\theta = 0.90$ pass rate
with $\alpha = 0.05$, $\beta = 0.10$, $\delta = 0.10$. The fixed-sample
test requires $n^* \approx 109$ trials (\cref{eq:sample-size}). Under
SPRT:
\begin{itemize}[leftmargin=*]
  \item If $p = 0.90$ (agent is fine): $\E[N] \approx 52$ trials
    (52\% savings).
  \item If $p = 0.80$ (agent has regressed): $\E[N] \approx 34$ trials
    (69\% savings).
\end{itemize}
The savings are most dramatic when the truth is far from the decision
boundary.
\end{example}

\subsection{Bayesian Regression Analysis}
\label{sec:stochastic:bayesian}

As a complement to the frequentist framework above, we also provide a
Bayesian perspective.

\begin{definition}[Bayesian Regression Probability]
\label{def:bayesian}
Given a Beta prior $p \sim \text{Beta}(a_0, b_0)$ on the pass rate,
baseline data $\mathbf{r}_b$ yielding posterior
$p_b \mid \mathbf{r}_b \sim \text{Beta}(a_0 + k_b, b_0 + n_b - k_b)$
where $k_b = \sum r_i^b$, and current data $\mathbf{r}_c$ yielding
posterior
$p_c \mid \mathbf{r}_c \sim \text{Beta}(a_0 + k_c, b_0 + n_c - k_c)$,
the \emph{Bayesian regression probability} is:
\begin{equation}
\label{eq:bayesian-regression}
  P_{\text{reg}} = \Pr[p_c < p_b - \delta \mid \mathbf{r}_b, \mathbf{r}_c]
\end{equation}
computed by sampling or numerical integration over the joint posterior.
The Bayesian verdict is $\FAIL$ if $P_{\text{reg}} > 1 - \alpha$,
$\PASS$ if $P_{\text{reg}} < \beta$, and $\INCONC$ otherwise.
\end{definition}

The Bayesian approach offers three advantages: it directly answers ``what
is the probability of regression?'' (rather than the indirect frequentist
framing), it naturally incorporates prior knowledge (e.g., the agent
historically passes 92\% of the time), and it provides a full posterior
distribution over the regression magnitude.

\subsection{Connection to Agent Behavioral Contracts}
\label{sec:stochastic:abc}

The \agentassert{} framework~\citep{bhardwaj2026abc} defines
\emph{$(p, \delta, k)$-satisfaction}: an agent $(p, \delta, k)$-satisfies
a behavioral contract $\phi$ if, over $k$ consecutive observations, the
empirical satisfaction rate exceeds $p$ and the behavioral drift metric
$D(t)$ remains below $\delta$.

\begin{proposition}[Verdict-Contract Correspondence]
\label{prop:verdict-contract}
Let $A$ be an agent with behavioral contract $\phi$ and associated
evaluator $E_\phi(x, o) = \mathbf{1}[o \models \phi]$. If the
stochastic test verdict $V(\mathbf{r}; \theta, \alpha) = \PASS$ with
threshold $\theta = p$ and $n \geq k$, then $A$
$(p', \delta, k)$-satisfies $\phi$ with
$p' \geq p - O(1/\sqrt{n})$ and probability at least $1 - \alpha$.
\end{proposition}

\begin{proof}[Proof sketch]
The \PASS{} verdict guarantees $\CI_{\text{lower}} \geq \theta = p$ at
level $1-\alpha$, which implies the true pass rate exceeds $p$. By
Hoeffding's inequality applied to $k$ consecutive observations within the
$n$ trials, the empirical rate over any window of $k$ observations
concentrates around the true rate. See \cref{app:proof:contract} for the
formal argument.
\end{proof}

This connection is powerful: it means that \agentassay{} test results
can directly certify contract compliance, closing the loop between
\emph{specification} (contracts), \emph{testing} (stochastic verdicts),
and \emph{enforcement} (runtime monitoring).

\subsection{Test Suite Semantics}
\label{sec:stochastic:suite}

\begin{definition}[Test Suite Verdict]
\label{def:suite}
A \emph{test suite} $\mathcal{S} = \{S_1, S_2, \ldots, S_m\}$ is a
collection of test scenarios. The \emph{suite verdict} aggregates
individual verdicts:
\begin{equation}
\label{eq:suite-verdict}
  V_{\text{suite}}(\mathcal{S}) =
  \begin{cases}
    \FAIL & \text{if } \exists\, j : V(S_j) = \FAIL \\
    \INCONC & \text{if } \nexists\, j : V(S_j) = \FAIL \text{ and }
      \exists\, j : V(S_j) = \INCONC \\
    \PASS & \text{if } \forall\, j : V(S_j) = \PASS
  \end{cases}
\end{equation}
\end{definition}

\begin{remark}[Multiple Testing Correction]
When running $m$ scenarios at significance level $\alpha$, the
family-wise error rate (FWER) can exceed $\alpha$. We apply the
Holm-Bonferroni correction: order $p$-values $p_{(1)} \leq \cdots \leq
p_{(m)}$ and reject $H_{0,(j)}$ if
$p_{(j)} \leq \alpha / (m - j + 1)$. This controls the FWER at level
$\alpha$ while preserving more power than Bonferroni correction.
\end{remark}
